% notations of definitions of unit 132 that this note uses and another note of the unit states
% (made by shared/mkspec.py from shared/notation-132.json; the same TeX as in the note that states them)
\newcommand{\RealPlaceSizes}[4]{\operatorname{RPS}(#1,#2,#3,#4)}
\newcommand{\haarRatio}[1]{\lambda(#1)}
\newcommand{\gauge}[2]{\Gamma(#1;#2)}
\newcommand{\HasFibreBound}[5]{\operatorname{Fib}_{#2,#4}(#1;#3)\le #5}
\newcommand{\gramSection}[1]{\gamma(#1)}
\newcommand{\invCorner}[1]{\operatorname{ic}(#1)}
\newcommand{\cornerDefect}[1]{\operatorname{cd}(#1)}
\newcommand{\IsBoundedInCell}[2]{#1\in\operatorname{Cell}_{\le #2}}
\newcommand{\diagLast}[1]{\operatorname{d}_{\varrho}(#1)}
\newcommand{\sectionCoeff}[2]{\mathsf{c}(#1,#2)}
\newcommand{\sectionFrame}[3]{\mathsf{f}(#1,#2,#3)}
\newcommand{\sectionMatrix}[3]{\mathsf{m}(#1,#2,#3)}
\newcommand{\Levels}[3]{\operatorname{Lev}(#1,#2,#3)}
\newcommand{\floorMeasure}[2]{\theta_{#2}(#1)}
\newcommand{\stackMeasure}[1]{\mu^{#1}}
\newcommand{\groupMeasure}[1]{\mathsf{M}(#1)}
\newcommand{\upperMeasure}[2]{\mathsf{M}_{\ge #2}(#1)}
\newcommand{\IsNear}[3]{#3\in\operatorname{Near}_{#2}(#1)}
\newcommand{\orthonormalPart}[1]{\mathsf{O}(#1)}
\newcommand{\closePairs}[1]{\operatorname{Cl}(#1)}
\newcommand{\floorValues}[4]{\operatorname{Val}_{#2}(#1;#3,#4)}
\newcommand{\volumeExponent}[1]{\mathsf{s}(#1)}
\newcommand{\floorWeight}[4]{\mathsf{w}_{#2}(#1;#3,#4)}
\newcommand{\floorModulus}[3]{\mathsf{d}_{#2}(#1;#3)}
\newcommand{\FloorMassLe}[4]{\operatorname{Mass}_{#2}(#1;#4)\le #3}
\newcommand{\upperPairing}[5]{\mathsf{Av}_{#2}(#1;#3)(#4,#5)}
\newcommand{\IsPairingKernel}[4]{(#3,#4)\in\operatorname{PK}_{#2}(#1)}
\newcommand{\invWeight}[3]{\mathsf{w}^{-1}_{#2}(#1;#3)}
\newcommand{\pivotModulus}[5]{\mathsf{Pv}_{#2}(#1;#3)(#4,#5)}
\newcommand{\weightedSquare}[5]{\mathsf{Sq}_{#2}(#1;#3,#4)(#5)}
\newcommand{\weightedSquareSum}[5]{\mathsf{Sm}_{#4}(#1;#2,#3)(#5)}
\newcommand{\selfNorm}[4]{\operatorname{SN}_{#2}(#1;#3,#4)}
\newcommand{\Gauges}[1]{\operatorname{Gau}(#1)}
\newcommand{\positionsOfSize}[2]{\operatorname{Pos}(#1;#2)}
\newcommand{\Sizes}[1]{\operatorname{Siz}(#1)}
\newcommand{\weylRatio}[3]{w_{#1}(#3)^{-1}w_{#1}(#2)}
\newcommand{\vecsOfSize}[1]{\mathrm{Vec}_{#1}}
\newcommand{\HasMatSize}[2]{\operatorname{sz}(#2)\le #1}
\newcommand{\IsHead}[2]{#2\in\operatorname{Hd}_{#1}}
\newcommand{\frameVolume}[1]{\Delta(#1)}
\newcommand{\frameXPart}[1]{\mathsf{x}_0(#1)}
\newcommand{\floorConst}[2]{C_{#2}(#1)}
\newcommand{\globalConst}[1]{C_{\mathrm{gl}}(#1)}
